\section{Ablations and robustness}
\label{sec:ablations}

This section varies one ingredient of PATTERN-DIF at a time and asks two questions: does the estimator remain valid, and how much efficiency is gained or lost relative to the default? The default is the method of the main study: strata $(G_i,K_i)$ with strata smaller than 50 merged within group, weighted generalised Neyman allocation \citep{neyman1934} computed with the ensemble-mean proxy $\pi_i$, a floor of 10 labels per stratum, the stratified difference estimator with $\lambda=1$ and the same proxy, and 95\% Wald intervals from the linearisation variance. The plain (Horvitz--Thompson) estimator, which uses the proxy only to allocate labels, is called HT or $\lambda=0$. Because the proxy enters both the allocation and the estimator, the proxy ablations report the two roles separately. Full details are in Supplementary Section~\SIref{si:ablations}.

\paragraph{Protocol.} We use the eleven cached pools and budgets $L\in\{250,500,1000\}$ with $L\le N/4$ (29 setting--budget cells), $R=150$ replications per cell and design, and the 108 comparison targets of the 36 pairs of nine models (targets that are exact for every design are excluded). Efficiency is the per-target ratio $\mathrm{MSE}_{\mathrm{SRS}}/\mathrm{MSE}_{\mathrm{design}}$ with SRS and the plain estimator as reference, summarised by the median over targets and then over cells. ``Rel.'' is the median over cells of $\mathrm{MSE}_{\mathrm{PATTERN\text{-}DIF}}/\mathrm{MSE}_{\mathrm{variant}}$ (above one favours the variant). With $R=150$ the minimum over 108 targets of coverage that is exactly nominal has mean 0.90, so worst-target coverage near 0.90 is consistent with nominal behaviour and values well below 0.88 are not.

\begin{table}[t]
\centering\small
\caption{Stratifier ablation (A1). Efficiency vs SRS (plain estimator): median over cells of the median over targets of $\mathrm{MSE}_{\mathrm{SRS}}/\mathrm{MSE}$. Rel.: median over the 29 cells of $\mathrm{MSE}_{\mathrm{PATTERN\text{-}DIF}}/\mathrm{MSE}_{\mathrm{variant}}$; Win: share of cells with Rel.\ $>1$. Coverage: mean over cells of the mean over targets (mean) and of the worst target (worst), nominal 0.95. Source: \texttt{results/label\_ablations\_dif.json}.}
\label{tab:abl-design}
\begin{tabular}{lccccccc}
\toprule
 & \multicolumn{3}{c}{Efficiency vs SRS} & & & \multicolumn{2}{c}{Coverage} \\
\cmidrule(lr){2-4}\cmidrule(lr){7-8}
Variant (estimator DIF unless stated) & $L=250$ & 500 & 1000 & Rel. & Win & mean & worst \\
\midrule
PATTERN-DIF (default) & 6.27 & 6.18 & 6.37 & 1.00 & -- & 0.945 & 0.877 \\
\addlinespace
\multicolumn{8}{l}{\emph{A1: stratifier}}\\
Group $\times$ agree/disagree & 4.33 & 4.65 & 4.81 & 0.74 & 0.07 & 0.946 & 0.881 \\
Group $\times$ pattern of 3 models & 2.82 & 2.98 & 3.43 & 0.38 & 0.00 & 0.941 & 0.867 \\
Group $\times K\times$ 2 proxy bins & 6.75 & 6.42 & 7.18 & 1.01 & 0.52 & 0.943 & 0.877 \\
Group $\times K\times$ 3 proxy bins & 6.79 & 5.90 & 6.61 & 0.99 & 0.38 & 0.945 & 0.890 \\
PATTERN, residual allocation & 5.64 & 6.02 & 6.68 & 0.96 & 0.45 & 0.945 & 0.886 \\
\addlinespace
\multicolumn{8}{l}{\emph{References}}\\
GROUP-DIF (strata = group) & 2.10 & 2.04 & 2.07 & 0.28 & 0.00 & 0.941 & 0.850 \\
SRS-DIF (PPI, $\lambda=1$) & 1.83 & 1.91 & 1.92 & 0.28 & 0.00 & 0.938 & 0.844 \\
Proportional allocation (no proxy) & 3.85 & 3.75 & 3.50 & 0.57 & 0.03 & 0.943 & 0.877 \\
Gap-targeted allocation (A5) & 5.36 & 5.26 & 5.22 & 0.74 & 0.00 & 0.941 & 0.866 \\
PATTERN, plug-in (HT) & 3.63 & 3.93 & 4.34 & 0.65 & 0.00 & 0.944 & 0.858 \\
\bottomrule
\end{tabular}

\end{table}

\paragraph{Stratifier.} On the same samples the difference estimator is more efficient than HT in all 29 cells (median ratio 1.67), and its efficiency against SRS is 6.3 at every budget, compared with 3.6 to 4.3 for HT. The gains are not additive: SRS with the difference estimator gains a factor 1.9, PATTERN with HT 4.0, and the product (7.5) exceeds the combined gain. The count $K_i$ is what matters (Table~\ref{tab:abl-design}): group $\times$ agree/disagree keeps 74\% of the default's precision, the joint pattern of the three models agreeing most with the ensemble keeps 38\%, and splitting $(G,K)$ cells by proxy bins is a tie (Rel.\ 0.99 to 1.01), whose extra strata eat the budget through the floor. A residual-variance allocation is not better (Rel.\ 0.96). The references that drop an ingredient lose clearly: GROUP-DIF and SRS-DIF keep 28\% and proportional allocation 57\%. Coverage hardly depends on the stratifier (mean 0.938 to 0.946).

\paragraph{Proxy.} The allocation is insensitive to the proxy; the estimator is not. With HT, every allocation proxy gives a median efficiency of 3.95 to 4.05 against 3.97 for the ensemble mean, including the reversed proxy $1-\pi$, a shuffled proxy and the constant 0.5, because the influence values vary within a stratum mainly through the differences $\hat y_{mi}-\hat y_{ji}$, which need no labels; what matters is that allocation is Neyman-type rather than proportional (Rel.\ 0.57). With the ensemble proxy kept in the estimator, alternative allocation proxies keep 99\% to 101\% of the default's precision. When the proxy is also used in the estimator, replacing the ensemble mean by one logistic regression costs 13\%, naive Bayes 42\%, the constant 16\%, a shuffled proxy 47\%, a noisy proxy 39\% and the reversed proxy 34\%. The difference estimator stays unbiased and its coverage is unaffected (mean 0.940 to 0.949), so a poor proxy costs efficiency only, but a corrupted proxy also removes the gain over HT (median DIF/HT ratio 0.96 for shuffled and noisy proxies against 1.67 for the ensemble). The per-proxy table and figure are in Supplementary Section~\SIref{si:ablations} (Supplementary Table~\SIref{tab:abl-proxy}).

\paragraph{The weight $\lambda$.} A $\lambda$ tuned per stratum and target and clipped to $[0,1]$ \citep[as in PPI++ and StratPPI;][]{angelopoulos2026ppipp,fisch2024stratppi} is slightly more efficient than $\lambda=1$ with the ensemble proxy (median Rel.\ 1.016) but undercovers: mean coverage 0.923 and worst target 0.799 on average (0.227 in Law school at $L=250$), against 0.945, 0.877 and 0.633 for $\lambda=1$. The weights are estimated from a few labels per stratum and the variance estimate ignores that. Neither alternative beats $\lambda=1$ in both efficiency and validity with a good proxy, so we keep $\lambda=1$; tuning is worth considering only when the proxy is known to be poor and some undercoverage is acceptable.

\paragraph{Number of models.} For $M\in\{2,3,5\}$ (random subsets of the nine models) the gain is governed by the share of the pool on which the compared models disagree rather than by $M$: the Spearman correlation between log efficiency and disagreement share is $-0.72$, $-0.74$ and $-0.76$ at $L=250$, 500 and 1000. At $L=250$ the median efficiency of PATTERN-DIF is 17.2 when the models disagree on less than 10\% of the pool, 12.5 for 10 to 20\%, 4.8 for 20 to 40\% and 3.8 above 40\%. The gain of DIF over HT grows with $M$, from 1.15 at $M=2$ to 1.67 at $M=9$. Coverage is not affected by $M$ (mean 0.945 to 0.950); see Supplementary Table~\SIref{tab:abl-models} and Supplementary Figure~\SIref{fig:abl-models}.

\begin{table}[t]
\centering\footnotesize
\setlength{\tabcolsep}{4pt}
\caption{Small-sample inference for PATTERN-DIF (A4). Mean: mean over cells and targets; worst: mean over cells of the worst-target coverage; min: minimum over cells of the worst-target coverage. Small pools: Churn, Spambase, Phoneme, Law school. Width: average interval width relative to Wald over all 29 cells. Floors: 'zero' replaces a zero stratum variance by the proxy residual variance, 'max' uses the larger of the two in every stratum.}
\label{tab:abl-inference}
\begin{tabular}{lcccccccc}
\toprule
 & \multicolumn{3}{c}{$L=250$, all} & \multicolumn{2}{c}{$L=250$, small} & \multicolumn{2}{c}{$L\ge 500$} & \\
\cmidrule(lr){2-4}\cmidrule(lr){5-6}\cmidrule(lr){7-8}
Method & mean & worst & min & worst & min & worst & min & Width \\
\midrule
Wald (default) & 0.941 & 0.852 & 0.633 & 0.827 & 0.633 & 0.892 & 0.833 & 1.00 \\
Wald + floor 'zero' & 0.947 & 0.876 & 0.780 & 0.893 & 0.873 & 0.893 & 0.833 & 1.00 \\
Wald + floor 'max' & 0.985 & 0.948 & 0.920 & 0.953 & 0.933 & 0.939 & 0.920 & 1.67 \\
$t$, Satterthwaite df & 0.945 & 0.865 & 0.700 & 0.850 & 0.700 & 0.893 & 0.833 & 1.01 \\
$t$ + floor 'zero' & 0.950 & 0.885 & 0.787 & 0.905 & 0.887 & 0.894 & 0.833 & 1.01 \\
Stratified bootstrap & 0.950 & 0.878 & 0.767 & 0.893 & 0.767 & 0.918 & 0.860 & 1.11 \\
\addlinespace
\multicolumn{9}{l}{\emph{References}} \\
SRS-DIF (PPI), Wald & 0.934 & 0.818 & 0.267 & 0.725 & 0.267 & 0.860 & 0.400 & 1.00 \\
SRS-DIF, Wald + floor 'max' & 0.976 & 0.929 & 0.900 & 0.942 & 0.900 & 0.929 & 0.880 & 1.23 \\
SRS-DIF, bootstrap & 0.924 & 0.741 & 0.100 & 0.678 & 0.100 & 0.818 & 0.153 & 1.12 \\
PATTERN plug-in (HT), Wald & 0.940 & 0.807 & 0.473 & 0.728 & 0.473 & 0.890 & 0.800 & -- \\
\bottomrule
\end{tabular}

\end{table}

\paragraph{Small-sample inference.} With the plain estimator, worst-target coverage was near 0.5 at $L=250$ in Churn and Law school because some strata showed no variation in the influence values. The difference estimator removes most of this: its residual $(Y_i-\pi_i)(u_i-v_i)$ varies even when no positive is sampled, no (replication, target) pair had a standard error of exactly zero, and the worst-target Wald coverage in Churn at $L=250$ rises from 0.47 to 0.91 (Table~\ref{tab:abl-inference}). Shortfalls remain: the worst-target Wald coverage is below 0.85 in five of 29 cells (Law school 0.63, MAGIC 0.78 and Adult-sex 0.82 at $L=250$, Bank 0.83 and Phoneme 0.83 at $L=500$). The mean standard error is 0.90 of the Monte Carlo standard deviation and the bias is small, so the intervals are too short rather than off-centre, consistent with a linearisation error in ratios with few sampled positives. A Satterthwaite $t$ quantile adds 1\% to the width and barely helps, and the stratified bootstrap is 11\% wider with worst target 0.878 and minimum 0.767. The proxy variance floor replaces a stratum variance that is exactly zero by its proxy residual variance ('zero') or takes the larger of the two in every stratum ('max'). The 'zero' floor leaves the width unchanged and lifts the minimum over cells from 0.633 to 0.780, entirely through Law school; only the 'max' floor reaches 0.92 in every cell, at 1.25 times the Wald width in the median (4.9 in Law school). The floors were designed after seeing the failure of the plain estimator on these pools, so the 'zero' result for Law school is in-sample; an independent check on synthetic pools (\texttt{results/floor\_check.json}; Supplementary Section~\SIref{sec:simulation}) found that the floor leaves the interval width unchanged (median ratio 1.00).

\paragraph{Ranking models by the equalised-odds gap.} Practitioners also rank models by their equalised-odds gap $EO_m=\max(|\Delta^{TPR}_m|,|\Delta^{FPR}_m|)$ \citep{hardt2016}. The plug-in estimate is biased upwards, as a maximum of absolute values must be. PATTERN-DIF is more accurate than SRS in the median (efficiency 1.94 to 2.21) and identifies the best model more often (mean probability 0.50 against 0.35 at $L=250$), but the difference estimator adds little over HT for this purpose, and PATTERN-DIF is less accurate than SRS in all Electricity cells (0.57 to 0.73), because the allocation spends labels on disagreement whereas the gap of one model depends on all its positive predictions. A variant whose allocation targets the $2M$ per-model gaps is the best design here (efficiency 2.30 to 2.61) but keeps only 74\% of the precision for the comparison contrasts. When the true margin to the next model is 0.002 or less, identification is hard for every design (Supplementary Table~\SIref{tab:abl-eo}).

\paragraph{Recommendations.} The count $K_i$ is the stratifier worth having, Neyman rather than proportional allocation matters, the allocation is insensitive to the proxy, and the gain depends on how often the compared models disagree. The proxy matters in the estimator, where a poor proxy costs up to half of the efficiency and removes the advantage over HT, although coverage is unaffected; the ensemble mean was never beaten there. Neither tuned $\lambda$ nor proxy-bin strata beat the default consistently, so we keep $\lambda=1$ and $K_i$ strata. We recommend reporting the Wald interval with the 'zero' floor, which costs nothing in width; the remaining undercoverage of rate contrasts at $L=250$ is not fixed by $t$ quantiles or the bootstrap, and only the conservative 'max' floor reaches nominal worst-target coverage. For ranking models by $EO_m$ itself, the gap-targeted allocation is better than the default.
